\documentclass[11pt,a4paper]{article}
\usepackage[T1]{fontenc}
\usepackage{lmodern,amsmath,amssymb,amsthm,mathtools,microtype}
\usepackage[margin=27mm]{geometry}
\usepackage[colorlinks=true,linkcolor=blue,citecolor=blue,urlcolor=blue]{hyperref}
\newtheorem{theorem}{Theorem}[section]
\newtheorem{proposition}[theorem]{Proposition}
\newtheorem{lemma}[theorem]{Lemma}
\newtheorem{corollary}[theorem]{Corollary}
\newcommand{\Q}{\mathbb Q}
\newcommand{\Z}{\mathbb Z}
\setlength{\emergencystretch}{2em}
\title{Constraints on Combining Pole Certificates}
\author{}\date{9 November 2026}
\begin{document}\maketitle
\begin{abstract}
Two elementary obstructions distinguish a joint pole certificate from
two independent certificates. First, a common nondegenerate conic
parameter forces the two inputs to have the same rational squareclass.
Two explicit prepared inputs admit no joint certificate with a common
parameter, including the degenerate values. Second, reflection-symmetric
contact sets admit an unbounded family on which the product of two
scaled contact functions is a square. We also identify the prime
divisors responsible for overlap between two scaled contact sets.
These statements concern the specified combinations; they are not
lower bounds for arbitrary Diophantine definitions.
\end{abstract}

\section{A common parameter in the pole conics}

Call a rational number $b$ \emph{prepared} if its reduced denominator
is a square and
\[
 \begin{gathered}
 0<b<17,\qquad
 b\in\Q_2^{\times2}\cap\Q_3^{\times2}\cap\Q_{17}^{\times2},\\
 v_2(b)\ge3,\qquad v_3(b)\ge1,\qquad v_{17}(b)\ge2.
 \end{gathered}
\]
For $B=17b$ and $V=B(B-b^2)$, the pole equations are
\begin{equation}\label{eq:pole}
 p^2=B(1-e)(1+3e),\qquad
 Bq^2=(s^2-eV)(s^2+3eV).
\end{equation}
All witnesses in this section are rational; zero witnesses are allowed.
The independent realization of these equations on prepared inputs is
the arithmetic input stated in~\cite{Dependencies}.

\begin{lemma}\label{lem:squareclass}
Let $b_0,b_1\in\Q^\times$. If both pole systems have a solution
with the same parameter $e\notin\{1,-1/3\}$, then $b_0/b_1$ is
a rational square.
\end{lemma}
\begin{proof}
Put $B_i=17b_i$. The first equations give
$p_i^2=B_i(1-e)(1+3e)\ne0$. Their quotient is
$b_0/b_1=(p_0/p_1)^2$.
\end{proof}

\begin{theorem}\label{thm:shared-e}
The prepared numbers
\[
 b_0=\left(\frac{204}{55}\right)^2,\qquad
 b_1=409\left(\frac{204}{1001}\right)^2
\]
have no pair of pole certificates with a common value of $e$.
\end{theorem}
\begin{proof}
Both reduced denominators are squares and both numbers lie in $(0,17)$.
Their valuations at $2,3,17$ are respectively $4,2,2$.
The additional factor $409$ is $1$ modulo $8$, $3$, and $17$,
so it is a square in all three local fields. Thus both numbers
are prepared. But
\[
                  \frac{b_1}{b_0}
                   =409\left(\frac{55}{1001}\right)^2
\]
is not a rational square. Lemma~\ref{lem:squareclass} therefore
forces $e=1$ or $e=-1/3$ in any proposed joint certificate.
For $b_0$, neither value permits the second pole equation:
this is the exact rank-certified obstruction proved in
\cite[Theorem~1.2 and Sections~2--3]{Fixed}.
That obstruction permits zero values of $q$ and hence applies
to the present convention. Both remaining values of $e$ are excluded.
\end{proof}

This theorem does not identify the parameters of the rational conic
charts. Indeed
\[
 e_i=\frac{\tau^2-B_i}{\tau^2+3B_i}
\]
usually gives different $e_i$ for the same $\tau$.
Thus a joint realization with a common chart parameter $\tau$
is a separate arithmetic question. Zariski density of the two
independent witness sets, proved in~\cite{Density}, does not
assert that their product meets the subvariety on which two
parameters are equal.

\section{A square product from symmetric contacts}

For $r\in\Q^\times$, let
\[
 \operatorname{Odd}(r)=\prod_{p:\,v_p(r)\ {\rm odd}}p.
\]
The product includes both positive and negative odd valuations;
it is finite. A nonzero rational square has $\operatorname{Odd}=1$.

Consider the eight contacts
\[
 \Theta=\{-3,-2,-1,0,2,3,4,5\}
        =\{1\pm1,1\pm2,1\pm3,1\pm4\}.
\]
Define
\[
 F(t)=\prod_{j=1}^4\bigl((t-1)^2-j^2\bigr),\qquad
 R(t)=\frac{F(t)}{(t-1)^{10}}.
\]
Both functions are invariant under $t\mapsto2-t$.

\begin{theorem}\label{thm:symmetric}
For every rational $a>1$, put
\[
             \lambda=\frac2{a+1},\qquad
             t_0=\lambda,\qquad t_1=a\lambda.
\]
Both $R(t_i)$ are defined, nonzero, and positive, and
\[
              R(t_0)=R(t_1),\qquad
              \operatorname{Odd}\bigl(R(t_0)R(t_1)\bigr)=1.
\]
Consequently no constants $C,C_0$ can make
\[
       h(a)\le C\log\operatorname{Odd}
                    \bigl(R(\lambda)R(a\lambda)\bigr)+C_0
\]
hold for all admissible positive rational $a,\lambda$.
The conclusion remains true if $a$ is restricted to
$a=\beta^4+T^4$ for any one fixed rational $T$ with $|T|>1$
and arbitrary rational $\beta$.
\end{theorem}
\begin{proof}
The identity $t_0+t_1=2$ gives $R(t_0)=R(t_1)$.
Since $0<t_0<1<t_1<2$, neither argument is a contact or
the pole $1$. At either argument $(t_i-1)^2<1$.
All four factors of $F(t_i)$ are therefore negative,
so their product is positive; the denominator is positive.
The product of the two equal values is a nonzero rational square.

Letting $a$ range through positive integers tending to infinity
contradicts the proposed height inequality. For the stated
restriction, fix $T$ and let $\beta$ range through positive
integers. Then $a>1$ and $h(a)\ge\log a$ tends to infinity,
whereas the odd radical remains one.
\end{proof}

More generally, if a rational contact function satisfies
$R(2c-t)=R(t)$ with $c\ne0$, the substitution
$\lambda=2c/(a+1)$ makes the two values equal wherever both
are defined. An unbounded admissible set of $a$ then gives
the same obstruction.

Theorem~\ref{thm:symmetric} concerns the product
$R(\lambda)R(a\lambda)$. Multiplication or division by a
nonzero rational square does not alter its conclusion.
Thus the odd-valuation argument supplies no information on this
family if that product is used as the squareclass of one pole
input. The theorem does not assert that all the other equations
of a proposed definition are soluble on that family. It also
does not apply to the two independent pole certificates.

\section{Where common contact primes occur}

Let $\Theta\subset\Q$ be finite. Choose a finite set $S$ of primes
containing the prime divisors of the nonzero contact numerators,
their denominators, and the nonzero differences of two contacts.
For $t\in\Q\setminus\Theta$, let $M_\Theta(t)$ be the product
of primes $p\notin S$ such that $v_p(t-\theta)>0$ is odd
for some $\theta\in\Theta$.

\begin{proposition}\label{prop:collisions}
Let $a,\lambda\in\Q^\times$, let $p\notin S$, and suppose that
$v_p(a)=0$. If $p$ belongs to both $M_\Theta(\lambda)$ and
$M_\Theta(a\lambda)$, then precisely one of the following occurs:
\begin{enumerate}
\item both contacts are zero, and $0\in\Theta$;
\item the contacts are nonzero $\theta,\eta\in\Theta$, and
      $v_p(a\theta-\eta)>0$.
\end{enumerate}
In the first case the two contact valuations are equal.
\end{proposition}
\begin{proof}
At each of the two arguments there is exactly one contact modulo
$p$, by the choice of $S$. Write them as $\theta,\eta$.
Then $\lambda\equiv\theta$ and $a\lambda\equiv\eta\pmod p$.
If $\theta=0$, the unit condition on $a$ forces $\eta=0$;
the converse follows in the same way. Their valuations agree
because $a$ is a unit. Otherwise both contacts are units, and
subtracting gives $a\theta-\eta\equiv0\pmod p$.
\end{proof}

In particular, a zero contact shared by the two scales disappears
from the odd radical of the product whenever $a$ is a square:
the corresponding factors multiply to $a\lambda^2$.
The other possible overlaps are governed by the fixed ratio set
\[
 {\cal R}=\{\eta/\theta:\theta,\eta\in\Theta\setminus\{0\}\}.
\]
If $a\notin{\cal R}$ and $a=A/B$ is reduced, every prime in
the second case divides the numerator of $a-r$ for some
$r\in{\cal R}$. Consequently their product $C(a)$ satisfies
\[
                  \log C(a)\le |{\cal R}|h(a)+C_\Theta.
\]
Indeed, write $r=u_r/v_r$ in lowest terms. Each relevant prime
divides $v_rA-u_rB$, whose absolute value is at most
$(|v_r|+|u_r|)\exp(h(a))$, and multiply these bounds.
This estimate explains the additional height term in a direct
overlap argument. It does not bound that term by a constant
independent of $a$, nor does it exclude a sharper argument
using further arithmetic information.

\section*{Exact verification}

The accompanying program in \nolinkurl{computation/pole_parameters},
\begin{center}\small
\nolinkurl{check_pole_merger_obstructions.py},
\end{center}
checks the prepared rational inputs, their nonsquare ratio,
the reflection identity coefficient by coefficient, and the
scaled rational-function identity. Its certificate records the
precise scope of these checks. The elliptic rank certificate
used in Theorem~\ref{thm:shared-e} remains the independently
reproducible certificate accompanying~\cite{Fixed}.

\begin{thebibliography}{9}
\bibitem{Dependencies}
\emph{Arithmetic dependencies of the construction with 29 variables},
companion paper, 8 November 2026,
\nolinkurl{papers/arithmetic_dependencies}.
\bibitem{Fixed}
\emph{Obstructions to fixed parameters in pole certificates},
companion paper, 9 November 2026,
\nolinkurl{fixed_pole_parameters_20261109_v16.tex}.
\bibitem{Density}
\emph{Zariski density of the prepared pole witnesses},
companion paper, 9 November 2026,
\nolinkurl{pole_witness_density_20261109_v14.tex}.
\end{thebibliography}
\end{document}
